\section*{Bab \ref{ch:b3:sequence-stratigraphy} --- Stratigrafi Sekuen}

\begin{solution}{exo:b3:sequence-stratigraphy:1}
$A = 0.5 + 0.1 = \qty{0.6}{m/ka}$; dalam \qty{20}{ka} terdapat akomodasi
sebesar \qty{12}{m}. Dengan muka laut turun, $A = -0.3 + 0.1 = \qty{-0.2}{m/ka}$:
penurunan relatif yang diperlambat oleh penurunan dasar cekungan.
\end{solution}

\begin{solution}{exo:b3:sequence-stratigraphy:2}
(a) \omterm{def:b3:sequence-stratigraphy:tract}{Traktus sistem tahap turun} (\omterm{def:b3:sequence-stratigraphy:regression}{regresi paksa}). (b) \omterm{def:b3:sequence-stratigraphy:tract}{Traktus sistem transgresif}. (c)
\omterm{def:b3:sequence-stratigraphy:tract}{Traktus sistem permukaan laut tinggi}. (d) \omterm{def:b3:sequence-stratigraphy:tract}{Traktus sistem permukaan laut rendah} (\omterm{def:b3:sequence-stratigraphy:regression}{regresi normal} setelah penurunan).
\end{solution}

\begin{solution}{exo:b3:sequence-stratigraphy:3}
Setiap \omterm{def:b3:sequence-stratigraphy:parasequence}{parasekuen} berakhir lebih dangkal daripada yang di bawahnya: susunan progradasional,
garis pantai bergerak ke arah laut, seperti pada \omterm{def:b3:sequence-stratigraphy:tract}{traktus sistem permukaan laut tinggi} atau rendah. Kebalikannya,
yakni setiap \omterm{def:b3:sequence-stratigraphy:parasequence}{parasekuen} berakhir lebih halus dan lebih dalam, bersifat
retrogradasional: \omterm{def:b3:sequence-stratigraphy:tract}{traktus sistem transgresif}.
\end{solution}

\begin{solution}{exo:b3:sequence-stratigraphy:4}
Lima pasangan lapisan per berkas ialah rasio eksentrisitas pendek (sekitar
\qty{100}{ka}) terhadap presesi (sekitar \qty{20}{ka}): setiap pasangan lapisan
merupakan satu siklus presesi, setiap berkas satu siklus eksentrisitas. Lajunya
adalah $1.6/100 = \qty{1.6}{cm/ka}$, dan satu pasangan lapisan setebal \qty{0.32}{m} berlangsung sekitar
\qty{20}{ka}, konsisten.
\end{solution}

\begin{solution}{exo:b3:sequence-stratigraphy:5}
$\mathrm dx_s/\mathrm dt = (4000 - 0.15\times20\,000)/50 = \qty{20}{m/ka}$,
ke arah laut. Garis pantai diam ketika $AW = Q$: $A = 4000/20\,000 =
\qty{0.2}{m/ka}$.
\end{solution}

\begin{solution}{exo:b3:sequence-stratigraphy:6}
$\mathrm dx_s/\mathrm dt = 4000/50 + 0.2/0.01 = 80 + 20 = \qty{100}{m/ka}$;
Penurunan muka laut saja menyumbang \qty{20}{m/ka}, atau seperlima.
\end{solution}

\begin{solution}{exo:b3:sequence-stratigraphy:7}
Garis pantai mundur ketika $Q < AW = 0.4\times25\,000 = \qty{10000}{m^3/ka}$ per
meter. Dengan $Q = 6000$: $(6000 - 10\,000)/10 = \qty{-400}{m/ka}$, artinya mundur
sejauh \qty{0.4}{km} setiap seribu tahun. Muka pantai yang rendah menghasilkan
transgresi yang cepat.
\end{solution}

\begin{solution}{exo:b3:sequence-stratigraphy:8}
$1/94.9 - 1/123.9 = \num{0.00247}\text{ siklus}/\unit{ka}$: rekaman yang lebih panjang daripada
$1/0.00247 = \qty{406}{ka}$ memisahkannya, dan \qty{600}{ka} (resolusi
0.00167) sudah cukup. $1/94.9 - 1/98.9 = 0.00042$: diperlukan sekitar
\qty{2400}{ka}, empat kali lebih panjang daripada rekaman.
\end{solution}

\begin{solution}{exo:b3:sequence-stratigraphy:9}
$16.4/(20\times40.98) = \qty{0.0200}{m/ka}$, \qty{2.0}{cm/ka}; lapisan abu
berjarak $3.1/0.020 = \qty{155}{ka}$.
\end{solution}

\begin{solution}{exo:b3:sequence-stratigraphy:10}
(a) $\tan\alpha = -0.4/200 = -0.002$: menurun sebesar \ang{0.11}. (b)
$0.2/50 = 0.004$: naik sebesar \ang{0.23}. Penampang seismik membentang puluhan
kilometer pada eksagerasi vertikal puluhan hingga ratusan, sehingga sudut itu tampak;
tebing sepanjang beberapa ratus meter hanya menunjukkan kenaikan satu meter atau kurang, yang
hilang dalam variasi lapisan.
\end{solution}

\begin{solution}{exo:b3:sequence-stratigraphy:11}
$1\,296\,000/54.725282 = \qty{23.68}{ka}$ ($g_5$); $/57.924383 =
\qty{22.37}{ka}$ ($g_2$); $/68.377912 = \qty{18.95}{ka}$ ($g_4$);
$/67.834313 = \qty{19.11}{ka}$ ($g_3$). $g_2 - g_5 = 3.199101''$ per tahun,
\qty{405.1}{ka}; $g_4 - g_5 = 13.65263''$, \qty{94.9}{ka}.
\end{solution}

\begin{solution}{exo:b3:sequence-stratigraphy:12}
Keduanya menghasilkan $0.3\times100 = \qty{30}{m}$ akomodasi pada laju yang sama; dengan
pasokan yang sama, geometrinya juga sama, karena sedimen hanya merespons terhadap
muka laut relatif. Perubahan eustatik berlaku sama di mana-mana:
sekuen yang sama dengan umur yang sama di tepi lain menunjuk ke riwayat pertama. Riwayat
penurunan yang disusun dari riwayat termal dan beban cekungan
(backstripping), atau catatan eustasi independen seperti isotop oksigen,
juga memisahkannya.
\end{solution}

\begin{solution}{pb:b3:sequence-stratigraphy:1}
\textbf{1.} $\dot E = 30\times2\pi/400\times\sin(2\pi t/400)$: paling besar
\qty{0.471}{m/ka}, naik pada \qty{100}{ka} dan \qty{500}{ka}, turun pada
\qty{300}{ka}.

\textbf{2.} \qty{0.020}{m/ka} pada engsel, 0.037 pada \qty{20}{km}, 0.062 pada
\qty{50}{km}, \qty{0.120}{m/ka} pada \qty{120}{km}.

\textbf{3.} $A < 0$ ketika $\sin(2\pi t/400) < -0.0346/0.471 = -0.0734$: $t =
200 + (400/2\pi)\arcsin0.0734 = \qty{204.7}{ka}$.

\textbf{4.} Dengan $\sigma = 0.0617$: $t = 400 - 63.66\arcsin0.131 =
\qty{391.6}{ka}$.

\textbf{5.} Sekitar \qty{187}{ka} dibandingkan dengan 200: eustasi yang turun harus lebih dulu
melampaui penurunan, dan berhenti mengalahkannya sebelum traktus permukaan laut rendah.

\textbf{6.} $0.045\times400 = \qty{18}{m}$, dibandingkan dengan rentang eustatik sebesar
\qty{60}{m}: sepertiga.

\textbf{7.} Ke arah darat dari 50 hingga sekitar \qty{120}{ka}, ke arah laut hingga sekitar
\qty{440}{ka}, ke arah darat hingga sekitar \qty{560}{ka}, ke arah laut lagi hingga
\qty{600}{ka} (\omterm{def:b3:seismology-imaging:ray-parameter}{titik balik} model berada pada 118.5, 439.5 dan
\qty{559}{ka}; tabel setiap \qty{50}{ka} hanya membatasinya).

\textbf{8.} $(6000 - 0.30\times15\,000)/40 = \qty{37.5}{m/ka}$, ke arah laut.

\textbf{9.} $6000/15\,000 = \qty{0.40}{m/ka}$.

\textbf{10.} $6000/35.7 + 0.424/0.012 = 168.1 + 35.3 = \qty{203}{m/ka}$,
berada dalam \qty{0.5}{\percent} dari model.

\textbf{11.} Pasokannya: \qty{83}{\percent} dari kecepatan.

\textbf{12.} $(-42.7 - 16.3)/(50.7 - 17.4) = \qty{-1.77}{m/km}$, \ang{0.10}.
Sejak saat itu, titik pertama telah mengalami penurunan $0.0345\times395.5 =
\qty{13.6}{m}$, titik kedua $0.0623\times208.5 = \qty{13.0}{m}$: hampir
sama, sehingga pada elevasi sekarang turunnya, \qty{59.0}{m}, berbeda dari
turun eustatik hanya sebesar \qty{0.6}{m}. Turun relatif di garis pantai pada
saat itu, eustasi dikurangi penurunan selama masa turun, lebih kecil.

\textbf{13.} 150 dan 200: permukaan laut tinggi; 250, 300, 350: tahap turun; 400:
permukaan laut rendah; 450, 500, 550: transgresif (garis pantai mundur tepat setelah
440); 600: permukaan laut tinggi.

\textbf{14.} 204.7 dan \qty{391.6}{ka} dibandingkan dengan 204.5 dan 391.5: sama,
dalam langkah model sebesar \qty{0.5}{ka}.

\textbf{15.} Di paparan, permukaan yang tersingkap selama masa turun: bagian atas
endapan permukaan laut tinggi dari engsel hingga \qty{17}{km}, berlanjut di atas
foreset tahap turun yang bertangga ke bawah; di cekungan, konformitas
korelatif di dasar atau puncak \omterm{def:b3:sequence-stratigraphy:tract}{traktus sistem tahap turun}, menurut konvensi. Paparan
sungguhan akan dipotong oleh sungai menjadi \omterm{def:b3:sequence-stratigraphy:condensed}{lembah torehan}, lapuk menjadi tanah,
dan sebagian tererosi.

\textbf{16.} $(18.5 - 52.8)/(559 - 439.5) = \qty{-0.29}{km/ka}$; \omterm{def:b3:sequence-stratigraphy:surfaces}{permukaan genang maksimum}
terbentuk di ujungnya.

\textbf{17.} Garis pantai bergerak cepat di atas dataran pesisir yang hampir datar:
pasokannya terperangkap di dekatnya, tersebar di wilayah yang luas, dan setiap tempat
menerima sedikit. Ke arah cekungan dari kaki lereng, melampaui sekitar \qty{55}{km}, dasar laut
kekurangan pasokan: \omterm{def:b3:sequence-stratigraphy:condensed}{penampang termampatkan} terbentuk di sana.

\textbf{18.} Resolusi $1/30 = \num{0.033}\text{ siklus}/\unit{m}$; Nyquist $1/(2\times
0.01) = \num{50}\text{ siklus}/\unit{m}$.

\textbf{19.} $7.5/2.0 = 3.75$, 1 dan $0.81/2.0 = 0.405$; rasio orbit
adalah $405/100 = 4.05$ dan $41/100 = 0.41$.

\textbf{20.} \qty{2.0}{m} untuk \qty{100}{ka}: \qty{2.0}{cm/ka}; $0.81/41 =
\qty{1.98}{cm/ka}$; $7.5/405 = \qty{1.85}{cm/ka}$. Panjang gelombang yang
dapat ditunjukkan \omterm{def:b3:sequence-stratigraphy:periodogram}{periodogram} dekat \qty{8}{m} adalah 10, 7.5 dan \qty{6}{m}: yang
sebenarnya \qty{8.1}{m} berada di antara dua nilai tersebut.

\textbf{21.} $30/0.020 = \qty{1500}{ka}$.

\textbf{22.} $23.7\times0.020 = \qty{0.474}{m}$. Jika diambil sampel setiap
\qty{25}{cm}, frekuensi Nyquist adalah 2 siklus per meter, panjang gelombang
\qty{0.5}{m}: presesi, pada 2.11 siklus per meter, berada di luar batas itu dan
terlipat kembali menjadi $4 - 2.11 = 1.89$, panjang gelombang palsu sebesar \qty{0.53}{m}.

\textbf{23.} Satu puncak tunggal cocok dengan suatu periode pada suatu laju, apa pun
asalnya; tiga puncak yang rasionya sama dengan rasio orbit tidak bergantung pada
laju, yang kemudian harus didapat sama dari masing-masing puncak.

\textbf{24.} Garis pantai bergerak dari \qty{17.37}{km} pada awal
turun hingga \qty{52.84}{km} pada akhir regresi permukaan laut rendah: pergeseran maksimum
ke arah cekungan sebesar \textbf{\qty[text-series-to-math=true]{35.5}{km}}.
\end{solution}
